% Generated draft: requires mathematical and source review.
\documentclass[11pt]{article}
\usepackage[margin=25mm]{geometry}
\usepackage{fontspec}
\setmainfont{Latin Modern Roman}
\usepackage{amsmath,amssymb,mathtools,unicode-math}
\setmathfont{Latin Modern Math}
\usepackage{enumitem,needspace,xurl,hyperref}
\setlength{\parindent}{0pt}
\setlength{\parskip}{5pt}
\setlength{\emergencystretch}{3em}
\newcommand{\N}{\mathbb{N}}
\newcommand{\Z}{\mathbb{Z}}
\newcommand{\Q}{\mathbb{Q}}
\newcommand{\R}{\mathbb{R}}
\newcommand{\C}{\mathbb{C}}
\newcommand{\F}{\mathbb{F}}
\newcommand{\A}{\mathbb{A}}
\newcommand{\PP}{\mathbb{P}}
\newcommand{\E}{\mathbb{E}}
\newcommand{\eps}{\varepsilon}
\newcommand{\dd}{\,\mathrm{d}}
\begin{document}
\section*{UnsolvedMath ID: 30000793}
\textbf{Source problem code:} catalogue entry
\subsection{Definitions and mathematical statement}
Fix a prime $p$. For a positive integer $D$, let $K_D=\Q(\sqrt{-D})$ and define
\[
r_p(D)=\dim_{\F_p}\bigl(\operatorname{Cl}(K_D)/p\operatorname{Cl}(K_D)\bigr),
\]
using additive notation for the ideal class group. The first group of questions asks, as $0<D<X$ and $X\to\infty$, for the mean of $p^{r_p(D)}$, the distribution of $r_p(D)$, and upper bounds depending on $D$.

If $\mathcal I(X)$ denotes the intended sample of those integers, the quantities to determine include
\[
\frac{1}{|\mathcal I(X)|}\sum_{D\in\mathcal I(X)}p^{r_p(D)},
\qquad
\frac{|\{D\in\mathcal I(X):r_p(D)=j\}|}{|\mathcal I(X)|}
\quad(j\geq0),
\]
including their limiting behavior where limits exist, and pointwise upper estimates for $r_p(D)$. The source does not say whether the sample comprises all positive integers, squarefree radicands, or a discriminant-restricted family. In particular $D$ as written is a radicand parameter, not automatically the absolute field discriminant; repeated radicands can represent the same quadratic field.

For the second group, retain $n(D)$ as the source's count of ``$S_5$-extensions of $\Q$ with squarefree discriminant $D$''. For the intended sample $\mathcal J(X)\subset\{D\in\Z:0<|D|<X\}$, the requested mean and distribution are
\[
\frac{1}{|\mathcal J(X)|}\sum_{D\in\mathcal J(X)}n(D),
\qquad
\frac{|\{D\in\mathcal J(X):n(D)=j\}|}{|\mathcal J(X)|}
\quad(j\geq0),
\]
together with pointwise upper bounds for $n(D)$ in terms of its discriminant size. It is not stated whether $S_5$ means a quintic field whose normal closure has Galois group $S_5$, or a Galois extension itself with group $S_5$; these are distinct degrees and discriminants. Counting up to $\Q$-isomorphism and whether nonsquarefree integers are excluded or assigned count zero also require confirmation. No predicted probability law, mean constant, or bound is invented.
\subsection{Short English statement}
Determine means, distributions and pointwise bounds for quadratic class-group $p$-ranks and, separately, counts of $S_5$ fields with a given squarefree discriminant. The sampling conventions and the meaning of $S_5$-extension are not specified.
\subsection{Sources}
\begin{itemize}
\item Supplied report: \emph{Explicit Methods in Number Theory} (2008), Oberwolfach Reports, DOI 10.4171/owr/2007/34, pp. 1957--2034; proposed year 2007. \url{https://doi.org/10.4171/owr/2007/34}.
\item Supplied Ellenberg--Venkatesh reference: \url{https://annals.math.princeton.edu/2006/163-2/p11}.
\item Local UnsolvedMath registry, permanent ID 30000793, catalogue entry; created 25 August 2026, updated 16 September 2026. Its corrected\_verified note reports comparison with the official report. The supplied 21 August 2026 literature triage labels the combined record open, explicitly calls its notation and averaging ranges incomplete, and reports important counting/upper-bound results without a single resolution of all the averages and distributions. This dated assessment is preserved without independently assigning a status to each target.
\end{itemize}
\end{document}
